\documentclass[11pt]{article}
\usepackage[utf8]{inputenc}
\usepackage[T1]{fontenc}
\usepackage{booktabs}
\usepackage{amsmath}
\usepackage{geometry}
\geometry{margin=2.5cm}

\title{Causal Fairness Analysis Report}
\author{Generated by FairMind and an LLM}
\date{2026-08-22}

\begin{document}
\maketitle

\section{Analysis setup}

\begin{tabular}{ll}
\toprule
Dataset & census\_income\_kdd \\
Protected attribute ($X$) & S1\_sex \\
Comparison & Female $\to$ Male \\
Outcome ($Y$) & T\_income\_level (1) \\
Mediator ($W$) & weeks\_worked\_in\_year \\
Confounder ($Z$) & education \\
\bottomrule
\end{tabular}

\section{Causal effects (FairMind, exact values)}

The five effects below are computed by FairMind through exact inference on the
fitted Bayesian network. These values are final and must not be recomputed or
altered.

\begin{tabular}{lr}
\toprule
Effect & Value \\
\midrule
Total Variation (TV) & 0.089620 \\
Total Effect (TE) & 0.081061 \\
Spurious Effect (SE) & 0.008559 \\
Direct Effect (DE) & 0.064497 \\
Indirect Effect (IE, decomposition form: TE = DE + IE) & 0.016563 \\
\bottomrule
\end{tabular}

\section{Qualitative interpretation}

\subsection{Total effect and spurious component (TV, TE, SE)}

The total disparity between the two groups is 0.089620, with a significant portion of this disparity being genuine, as the total effect (TE) is 0.081061. The spurious effect (SE) is 0.008559, indicating that a small part of the observed disparity is driven by the confounder (education) rather than the protected attribute (sex).

\subsection{Direct effect (DE)}

The direct effect (DE) is 0.064497, which suggests a substantial direct discrimination against the protected group (females). This direct effect is notably larger than the total effect, indicating that the direct pathway from the protected attribute to the outcome is a major contributor to the disparity.

\subsection{Indirect effect (IE)}

The indirect effect (IE) is 0.016563, which amplifies the direct effect since it has the same sign. This means that the mediator (weeks worked in year) adds to the disparity, making the total effect larger than the direct effect alone.

\section{Recap Questions}

\begin{enumerate}
\item Is there direct discrimination against the protected group? \textbf{LLM:} YES \quad \textbf{Expected:} YES \quad (correct)
\item Does the indirect effect through the mediator mitigate the total effect? \textbf{LLM:} NO \quad \textbf{Expected:} NO \quad (correct)
\item Does the observed total effect exceed the practical relevance threshold? \textbf{LLM:} YES \quad \textbf{Expected:} YES \quad (correct)
\item Does the spurious component (SE) contribute substantially to the observed difference? \textbf{LLM:} NO \quad \textbf{Expected:} NO \quad (correct)
\item Is the direct effect the dominant component compared to the indirect effect? \textbf{LLM:} YES \quad \textbf{Expected:} YES \quad (correct)
\end{enumerate}

\vspace{1em}
\noindent\footnotesize
The recap answers follow deterministic threshold rules applied to the values above: direct discrimination when $|\mathrm{DE}| > 0.05$; a mediator channel counted as negligible when $|\mathrm{IE}| < 0.005$; practical relevance when $|\mathrm{TV}| > 0.05$; a substantial spurious component when $|\mathrm{SE}| / |\mathrm{TV}| > 0.1$.


\vspace{1em}
\noindent\textbf{Answer consistency:} 5/5 (100.0\%). The expected answers are derived deterministically from the FairMind values alone, through threshold rules, and were added to the document after it was generated: the model never saw them.

\end{document}